\documentclass[10pt,a4paper]{amsart}
\usepackage[margin=19mm]{geometry}
\usepackage{amsmath,amssymb,mathtools}
\usepackage[scr=rsfs,cal=euler]{mathalfa}
\usepackage{xfrac}
\usepackage[unicode,hidelinks,bookmarks=false]{hyperref}
\newcommand{\ie}{i.e.\ }
\newcommand{\cf}{cf.\ }
\newcommand{\resp}{respectively\ }
\newcommand{\Subsection}{\subsection}
\providecommand{\nobreakdash}{\nobreak\hbox{-}}
\setlength{\emergencystretch}{2em}
\multlinegap0pt
\usepackage[capitalize]{cleveref}
\newcommand{\E}{{\mathbf E}}
\renewcommand{\P}{\mathbf P}
\renewcommand{\le}{\leqslant}
\renewcommand{\ge}{\geqslant}
\newcommand{\wh}{\widehat}
\newcommand{\di}{\mathrm{d}}
\newcommand{\M}{M}
\newcommand{\X}{\mathcal{X}}
\newcommand{\Qs}{\mathcal{Q}}
\newcommand{\kl}{\operatorname{KL}}
\newcommand{\kll}[2]{\kl ({#1}, {#2})}
\newcommand{\hel}{\operatorname{H}}
\newcommand{\hels}{\hel^2}
\newtheorem{lemma}{Lemma}
\newtheorem{proposition}{Proposition}
\title[Two-model Bayesian averaging]{Sharp deviation bounds for two-model Bayesian averaging}
\author[Compton et al.]{Spencer Compton, G\'{a}bor Lugosi, Jaouad Mourtada, Jian Qian, Nikita Zhivotovskiy}
\date{Source: arXiv:2602.16142v2; CC BY 4.0}
\begin{document}\maketitle
\noindent\textit{Counting scope.} The counted unit is Proposition \texttt{prop:bayes-two} and its complete author proof. The original definitions, KL--Hellinger comparison and almost-likelihood lemma with its proof are uncounted support. Mathematical excerpts are verbatim; labels are retained and standalone numbering is editorial. All five authors are credited, and the four references follow the fixed-version bibliography.\par\medskip
\noindent\textit{Source and licence.} Original author TeX, arXiv:2602.16142v2; CC BY 4.0, verified on the official versioned article page.\par\medskip
\section{Original definitions and comparison tools}
Our goal is to construct, based on $X_1,\ldots,X_n$, a density estimator $\widehat{p}$ for which the Kullback--Leibler divergence $\kl(p^\star,\widehat{p})$ is small with high probability (where we recall that $\kll{p}{q} = \int_\X p \log (p/q) \di \mu$ for two probability densities $p, q$ on $\X$ with respect to $\mu$).

Let $[k]$ denote $\{1, 2, \ldots, k\}$. $\wh L$ is the normalized negative log-likelihood, namely
\begin{equation}
  \label{eq:def-log-likelihood}
  \wh L(p) = -\frac{1}{n}\sum\limits_{i = 1}^n\log p(X_i).
\end{equation}
Squared Hellinger divergence is denoted as
\begin{equation}
  \label{eq:def-hellinger}
  \hels(p,q) \triangleq \frac{1}{2} \int_{\mathcal X} \big(\sqrt{p(x)}-\sqrt{q(x)} \big)^2\di \mu(x).
\end{equation}

\begin{lemma}
    \label{lem:kltohellinger}
Assume that $p,q$ are two densities on $\X$ with respect to $\mu$.
It holds that
\[
\kl(p,q) \le \frac{2}{(\sqrt{e}-1)^2}\,\hels(p,q)\,
\max\left\{1,\ \sup_{x\in\X}\log\left(\frac{p(x)}{q(x)}\right)\right\}.
\]
\end{lemma}

\noindent\textit{Standard prerequisite.} We use the preceding standard KL--Hellinger comparison in the general-density form stated by the authors, who cite \cite{birge1998minimum,birge1983approximation,tang2022divergence,birge1993rates}. Tang, Fact~4, states the finite-alphabet version with squared Hellinger distance defined without $1/2$; it supplies a constant and normalization comparison. The source appendix makes an incorrect global monotonicity claim, so that auxiliary proof is excluded from the dependency chain. It remains in the complete native source. The counted Bayesian argument below is unchanged.\par\medskip
\subsection*{Original likelihood-ratio lemma and proof}
\begin{lemma}
\label{lem:almostmle}
Let $\mathcal{P}=\{p_1,\ldots,p_\M\}$ be a finite set of densities with respect to a measure $\mu$ on $\mathcal{X}$.
Fix $p^\star\in\mathcal{P}$ and let $X_1,\ldots,X_n$ be i.i.d.\ distributed according to $p^\star$.
For $\delta \in (0, 1)$, set
\[
\wh \Qs_1
=\bigg\{p\in\mathcal{P}:\ \wh L(p)-\min_{q\in\mathcal{P}}\wh L(q)\le\frac{2\log(2\M/\delta)}{n}\bigg\}.
\]
Then, with probability at least $1-\delta$, one has $p^\star \in \wh \Qs_1$ and
\[
\max_{p \in \wh \Qs_1}\hel^2(p^\star,p) \le\frac{2\log(2\M/\delta)}{n}.
\]
\end{lemma}

\begin{proof}
Fix $t>0$. For any $p\in\mathcal{P}$, by Markov's inequality,
\begin{align*}
\P\left(\wh L(p)-\wh L(p^\star)\le t\right)
&=\P\left(\sqrt{\frac{\prod_{i=1}^n p(X_i)}{\prod_{i=1}^n p^\star(X_i)}}\ge \exp(-nt/2)\right)\\
&\le \exp(nt/2)\,\E\left[\sqrt{\frac{\prod_{i=1}^n p(X_i)}{\prod_{i=1}^n p^\star(X_i)}}\right]\\
&= \exp(nt/2)\,\big(1-\hel^2(p^\star,p)\big)^{n}
\le\exp\Big(\tfrac{nt}{2}-n\hel^2(p^\star,p)\Big),
\end{align*}
where we used $1-u\le \exp(-u)$. Let $\mathcal{P}_t=\{p\in\mathcal{P}:\hel^2(p^\star,p)\ge t\}$. By a union bound,
\[
\P\left(\exists p\in\mathcal{P}_t:\ \wh L(p)-\wh L(p^\star)\le t\right)
\le\sum_{p\in\mathcal{P}_t} \exp(-nt/2)
\le M\,\exp(-nt/2).
\]
Note that for any $p\in\mathcal{P}$,
$\wh L(p)-\min_{q\in\mathcal{P}}\wh L(q)\le t$ implies $\wh L(p)-\wh L(p^\star)\le t$,
since $\min_{q \in \mathcal{P}}\wh L(q)\le \wh L(p^\star)$. Therefore,
\[
\P\left(\exists p\in\mathcal{P}: \wh L(p)-\min_{q \in \mathcal{P}} \wh L(q)\le t\ \text{and}\ \hel^2(p^\star,p)\ge t\right) \le M \exp(-nt/2).
\]
For the second part, apply the same argument with $-t$ in place of $t$. For any $p\in\mathcal{P}$, it holds that
\begin{align*}
\P\left(\wh L(p)-\wh L(p^\star)\le -t\right)
&=\P\left(\sqrt{\frac{\prod_{i=1}^n p(X_i)}{\prod_{i=1}^n p^\star(X_i)}}\ge \exp(nt/2)\right)
\\
&\le \exp(-nt/2)\,\Big(\int\sqrt{p(x)p^\star(x)}\,d\mu(x)\Big)^{n}
\le \exp(-nt/2).
\end{align*}
By a union bound,
\[
\P\left(\min_{q\in\mathcal{P}} \wh L(q) \le \wh L(p^\star)-t\right)
\le \sum_{p\in\mathcal{P}} \P\left(\wh L(p)-\wh L(p^\star)\le -t\right)
\le M \exp(-nt/2).
\]
Equivalently, we have
\[
\P\left(\wh L(p^\star)-\min_{q\in\mathcal{P}} \wh L(q) \geq t\right)\le M \exp(-nt/2).
\]
Choose $t={2\log(2\M/\delta)}/{n}$. Then each of the two bad events above has probability at most $\delta/2$, so by a union bound they fail simultaneously with probability at most $\delta$. Hence, with probability at least $1-\delta$, every $p\in\wh \Qs_1$ satisfies $\hels(p^\star,p) \leq t$, and also $\wh L(p^\star)-\min_{q\in\mathcal{P}} \wh L(q)\le t$ hence $p^\star\in \wh \Qs_1$. This yields the claim.
\end{proof}

\section{Counted proposition and complete author proof}
\begin{proposition}
\label{prop:bayes-two}
Let $p^\star,q$ be two densities on $\X$ with respect to $\mu$, and let $X_1,\ldots,X_n$ be i.i.d.\ from $p^\star$. Consider Bayesian model averaging with a uniform prior, which outputs
\[
\widehat p=(1-\widehat\alpha)p^\star+\widehat\alpha\,q,\qquad
\widehat\alpha=\frac{\prod_{i=1}^n q(X_i)}{\prod_{i=1}^n p^\star(X_i)+\prod_{i=1}^n q(X_i)}.
\]
Then for any $\delta\in(0,1/2)$, with probability at least $1-\delta$,
\[
\kl(p^\star,\widehat p)\le \min\left\{\frac{10\log^2(2/\delta)}{n}, 5\log\left(\frac{2}{\delta}\right)\right\}.
\]
Moreover, taking $\X=\{0,1\}$, letting $q$ be the point mass at $1$, and letting $p^\star$ be the Bernoulli law with $(1-p^\star(0))^n=\delta$, we have, for any $\delta\in(0,e^{-2}]$ and with probability at least $\delta$,
\[
\kl(p^\star,\widehat p)\ge \min\left\{\frac{\log^2(1/\delta)}{4n},\, \frac{1}{4}\log\left(\frac{1}{\delta}\right)\right\}.
\]
\end{proposition}

\begin{proof}
By Markov's inequality and $\E\big[\prod_{i=1}^n q(X_i)/p^\star(X_i)\big]\le 1$,
\[
\P\left(\frac{\prod_{i=1}^n q(X_i)}{\prod_{i=1}^n p^\star(X_i)}\ge \frac{1}{\delta}\right)\le \delta.
\]
Hence, on an event of probability at least $1-\delta$ we have $\frac{\prod q(X_i)}{\prod p^\star(X_i)}<1/\delta$, which implies
\[
1-\widehat\alpha=\frac{1}{1+\frac{\prod q(X_i)}{\prod p^\star(X_i)}}\ge \frac{\delta}{1+\delta}\ge \frac{\delta}{2}.
\]
Therefore, for all $x$, 
\[
\widehat p(x)=(1-\widehat\alpha)p^\star(x)+\widehat\alpha q(x)\ge (1-\widehat\alpha)p^\star(x)\ge \frac{\delta}{2}p^\star(x),
\]
and so
\begin{equation}
\label{eq:env-bayes}
\sup_{x}\log\left(\frac{p^\star(x)}{\widehat p(x)}\right)\le \log\left(\frac{2}{\delta}\right).
\end{equation}
By Lemma~\ref{lem:kltohellinger},
\begin{equation}
\label{eq:hellingerbound}
\kl(p^\star,\widehat p)\le 5\hels(p^\star,\widehat p)\log\left(\frac{2}{\delta}\right).
\end{equation}
We consider the cases. First, if $2\log(2/\delta)/n > 1$, then the minimum in the statement is achieved at the second term and since $\hels(p^\star,\widehat p) \le 1$, we have
\[
\kl(p^\star,\widehat p)\le 5\log\left(\frac{2}{\delta}\right),
\]
which proves the claim. Otherwise, we have $2\log(2/\delta)/n \le 1$, and we again consider two cases. In the first case, assume $\hels(p^\star,q)\le \frac{2\log(2/\delta)}{n}$.
Using \eqref{eq:env-bayes} and the convexity of the squared Hellinger distance in its second argument, $\hels(p^\star,\widehat p)\le \hels(p^\star,q)$, which implies
\[
\kl(p^\star,\widehat p)\le \frac{10\log^2(2/\delta)}{n}.
\]
Otherwise, we have $\hels(p^\star,q) > \frac{2\log(2/\delta)}{n}$. In this case we want to upper bound $\widehat{\alpha}$. We have as in the proof of \Cref{lem:almostmle} for any $t > 0$,
\[
\P\left(
\frac{\prod_{i=1}^n q(X_i)}{\prod_{i=1}^n p^\star(X_i)} \ge t
\right)
\le
\frac{1}{\sqrt{t}}\bigl(1-\hels(p^\star,q)\bigr)^n
\le
\frac{1}{\sqrt{t}} \exp\bigl(-n\hels(p^\star,q)\bigr).
\]
We choose $t = \frac{2\log(2/\delta)}{n\hels(p^\star,q)}$. Note that $t<1<1/\delta$. By the above line we have
\[
\P\left(
\frac{\prod_{i=1}^n q(X_i)}{\prod_{i=1}^n p^\star(X_i)}
\ge
\frac{2\log(2/\delta)}{n\hels(p^\star,q)}
\right)
\le
\sqrt{\frac{n\hels(p^\star,q)}{2\log(2/\delta)}}\exp(-n\hels(p^\star,q))
\le
\exp(-2\log(2/\delta))
\le
\delta,
\]
where we used that $\sqrt{x}\exp(-x)$ is decreasing for $x \ge 1/2$ and $n\hels(p^\star,q) > 2\log(2/\delta) > 1/2$. Therefore, on the complementary event,
$
\widehat{\alpha}
\le
\frac{\prod_{i=1}^n q(X_i)}{\prod_{i=1}^n p^\star(X_i)}
<
\frac{2\log(2/\delta)}{n\hels(p^\star,q)},
$
and \eqref{eq:env-bayes} also holds since $t<1/\delta$. 
Hence, by \eqref{eq:hellingerbound} and the convexity of the squared Hellinger distance in its second argument, we have on this event
\[
\kl(p^\star,\widehat p)
\le
5\hels(p^\star,\widehat p)\log\left(\frac{2}{\delta}\right)
\le
5\widehat{\alpha}\hels(p^\star,q)\log\left(\frac{2}{\delta}\right)
\le
\frac{10\log^2(2/\delta)}{n}.
\]
The upper bound follows.

For the lower bound, let $\X=\{0,1\}$. Let $q$ be the point mass at $1$ (so $q(1)=1,q(0)=0$). Fix $p\in(0,1)$ such that $(1-p)^n=\delta$, and let $p^\star$ be the Bernoulli law with $p^\star(0)=p$ and $p^\star(1)=1-p$. On the event
\[
A=\{X_1=\cdots=X_n=1\},
\]
which satisfies $\P(A)=(1-p)^n=\delta$, we have
\[
\widehat\alpha=\frac{\prod_{i=1}^n q(X_i)}{\prod_{i=1}^n p^\star(X_i)+\prod_{i=1}^n q(X_i)}=\frac{1}{1+\delta},
\qquad
\widehat p(0)=(1-\widehat\alpha)p=\frac{p\delta}{1+\delta}.
\]
By the Bernoulli KL formula,
\[
\kl(p^\star,\widehat p)
=(1-p)\log\left(\frac{(1+\delta)(1-p)}{1+\delta-p\delta}\right)
+p\log\left(\frac{1+\delta}{\delta}\right).
\]
Since $1+\delta-p\delta\le 1+\delta$, the ratio $\frac{(1+\delta)(1-p)}{1+\delta-p\delta}$ lies in $[1-p,1]$. Hence
\[
(1-p)\log\left(\frac{(1+\delta)(1-p)}{1+\delta-p\delta}\right)\ge (1-p)\log(1-p)\ge -p,
\]
where the last inequality follows from $-\log(1-x)\le \frac{x}{1-x}$. Therefore, on $A$,
\begin{equation}
\label{eq:lb-core}
\kl(p^\star,\widehat p)\ge p\log\left(\frac{1+\delta}{\delta}\right)-p \ge p\log\left(\frac{1}{\delta}\right)-p.
\end{equation}
We now show that $p\ge \frac{\log(1/\delta)}{n+\log(1/\delta)}$.
Using $1-x\ge \exp\left(-\frac{x}{1-x}\right)$ for $x\in(0,1)$, we have
\[
\delta=(1-p)^n\ge \exp\left(-\frac{pn}{1-p}\right),
\]
which implies $p\ge \frac{\log(1/\delta)}{n+\log(1/\delta)}$.
Thus, if $\log(1/\delta)\le n$ then $p\ge \frac{1}{2}\frac{\log(1/\delta)}{n}$.
Since $\delta\le e^{-2}$ implies $\log(1/\delta)\ge 2$, from \eqref{eq:lb-core} we get
\[
\kl(p^\star,\widehat p)\ge p\bigl(\log(1/\delta)-1\bigr)\ge \frac{1}{2}\,p\log\left(\frac{1}{\delta}\right).
\]
Therefore, for $\delta\in(0,e^{-2}]$, it holds that
\[
\kl(p^\star,\widehat p)\ \ge\ \frac{1}{2}\,p\log\left(\frac{1}{\delta}\right)
\ \ge\ \begin{cases}
\displaystyle \frac{1}{4}\frac{\log^2(1/\delta)}{n}, & \text{if }\ \log(1/\delta)\le n,\\[1.2ex]
\displaystyle \frac{1}{4}\log\left(\frac{1}{\delta}\right), & \text{if }\ \log(1/\delta)\ge n,
\end{cases}
\]
where, in the second case, we also used $p\ge \frac{\log(1/\delta)}{n+\log(1/\delta)}\ge \frac{1}{2}$. Since $\P(A)=\delta$, these lower bounds hold with probability at least $\delta$. This completes the proof.
\end{proof}

\begin{thebibliography}{99}
\bibitem[BM98]{birge1998minimum} Lucien Birg{\'e} and Pascal Massart, Minimum contrast estimators on sieves: exponential bounds and rates of convergence, \emph{Bernoulli} \textbf{4} (1998), no.~3, 329--375.
\bibitem[Bir83]{birge1983approximation} Lucien Birg{\'e}, Approximation dans les espaces m{\'e}triques et th{\'e}orie de l'estimation, \emph{Zeitschrift f{\"u}r Wahrscheinlichkeitstheorie und verwandte Gebiete} \textbf{65} (1983), 181--237, Springer.
\bibitem[Tan22]{tang2022divergence} Jennifer Tang, \emph{Divergence Covering}, PhD thesis, Massachusetts Institute of Technology, 2022.
\bibitem[BM93]{birge1993rates} Lucien Birg{\'e} and Pascal Massart, Rates of convergence for minimum contrast estimators, \emph{Probability Theory and Related Fields} \textbf{97} (1993), 113--150, Springer.
\end{thebibliography}
\end{document}
